\section*{Chapitre \ref{ch:g9:uniform-varied-motion} --- Le mouvement~: uniforme et varié}

\begin{solution}{exo:g9:uniform-varied-motion:1}
Des marques des positions de l'objet à des instants également
espacés. Deux façons d'en obtenir~: parcourir une vidéo image par
image~; l'enregistreur à points qui frappe une bande de papier tirée
(ou une photographie au stroboscope).
\end{solution}

\begin{solution}{exo:g9:uniform-varied-motion:2}
Écarts égaux~: uniforme. $v = 0.04 \div 0.1 = \qty{0.4}{m/s}$.
\end{solution}

\begin{solution}{exo:g9:uniform-varied-motion:3}
Écarts croissants~: accéléré. Premier intervalle~:
$0.01 \div 0.1 = \qty{0.1}{m/s}$~; dernier~:
$0.07 \div 0.1 = \qty{0.7}{m/s}$.
\end{solution}

\begin{solution}{exo:g9:uniform-varied-motion:4}
Le radar lit la \omterm{def:g9:uniform-varied-motion:instant}{vitesse instantanée} --- l'allure à l'instant du
portail. Le «~vingt minutes de porte à porte~» calcule la moyenne ---
distance totale sur \omterm{def:g2:measuring-time:duration}{durée} totale, arrêts et pointes confondus.
\end{solution}

\begin{solution}{exo:g9:uniform-varied-motion:5}
Dans les \omterm{def:g6:describing-motion:uniform}{mouvements uniformes}~: l'allure ne change jamais, si bien que
tout intervalle --- long ou bref comme un clin d'œil --- redonne le
même nombre, et les deux notions s'effondrent en une seule.
\end{solution}

\begin{solution}{exo:g9:uniform-varied-motion:6}
D'environ \qty{9.8}{m/s} chaque seconde --- numériquement le $g$
local, \qty{9.8}{N/kg}~: la plus profonde rime de l'année.
\end{solution}

\begin{solution}{exo:g9:uniform-varied-motion:7}
Des \omterm{def:g2:pushes-and-pulls:force}{forces} compensées (ou absentes) laissent le mouvement inchangé~:
le repos reste repos, et le mouvement rectiligne régulier se poursuit.
Le corps libre de \omterm{def:g2:pushes-and-pulls:force}{forces} file à jamais~; l'égalité à quatre de l'avion
de ligne --- poussée contre traînée, portance contre \omterm{def:g4:weight-and-mass:weight}{poids} --- est
exactement pourquoi ses \qty{900}{km/h} se maintiennent.
\end{solution}

\begin{solution}{exo:g9:uniform-varied-motion:8}
La sonde, libre de \omterm{def:g2:pushes-and-pulls:force}{forces} dans le vide, file depuis des décennies
moteurs froids --- un mouvement qui n'a besoin d'aucune \omterm{def:g2:pushes-and-pulls:force}{force} pour
continuer. La bicyclette qui «~prouve~» le contraire ralentit sous les
frottements et l'\omterm{def:g2:air-around-us:air}{air}, des \omterm{def:g2:pushes-and-pulls:force}{forces} non compensées vers l'arrière et
déguisées~; sur la glace, le déguisement aminci, elle file loin.
\end{solution}

\begin{solution}{exo:g9:uniform-varied-motion:9}
Écarts décroissants~: ralenti. \omterm{def:g7:motion-average-speed:formula}{Vitesses}~: $0.6$, $0.5$, $0.4$, $0.3$,
\qty{0.2}{m/s} --- en baisse de \qty{0.1}{m/s} par intervalle. Si le
motif se maintient, deux intervalles de plus (\qty{1.0}{cm}, puis
rien) amènent la bande au repos.
\end{solution}

\begin{solution}{exo:g9:uniform-varied-motion:10}
\qty{0.1}{m/s} par \qty{0.1}{s}, c'est \qty{1}{m/s} gagné par seconde
--- environ le dixième des \qty{9.8}{m/s} par seconde de la chute
libre~: une poussée douce à côté de celle de la Terre.
\end{solution}

\begin{solution}{exo:g9:uniform-varied-motion:11}
Première phase~: le \omterm{def:g4:weight-and-mass:weight}{poids} l'emporte sur la jeune résistance de l'air
--- \omterm{def:g2:pushes-and-pulls:force}{force} non compensée vers le bas, la \omterm{def:g7:motion-average-speed:formula}{vitesse} croît. Voile ouverte~:
la traînée de l'\omterm{def:g2:air-around-us:air}{air} qui s'épanouit monte jusqu'à égaler exactement le
\omterm{def:g4:weight-and-mass:weight}{poids} --- \omterm{def:g2:pushes-and-pulls:force}{forces} compensées, et la chute s'installe aux \qty{5}{m/s}
constants que promet la phrase.
\end{solution}

\begin{solution}{exo:g9:uniform-varied-motion:12}
Enregistrement~: filmer le lâcher et parcourir les images toutes les
\qty{0.1}{s} (ou attacher la voiture à l'enregistreur à points)~;
\omterm{def:g6:measurement-in-science:measuring}{mesurer} écart après écart, sur la rampe puis sur le tapis. Nombres~:
les \omterm{def:g7:motion-average-speed:formula}{vitesses} d'intervalle, listées dans l'ordre. Verdict attendu~:
acte un, écarts croissants --- accéléré sur la rampe (la part non
compensée du \omterm{def:g4:weight-and-mass:weight}{poids} le long de la pente)~; acte deux, écarts
décroissants --- ralenti sur le tapis (la \omterm{def:g2:pushes-and-pulls:force}{force} non compensée du
frottement, vers l'arrière). Un seul jouet, les deux visages de la
phrase tranquille.
\end{solution}

\begin{solution}{pb:g9:uniform-varied-motion:1}
\textbf{1.} Uniforme~; $v = 0.08 \div 0.1 = \qty{0.8}{m/s} \approx
\qty{2.9}{km/h}$ --- une promenade.
\textbf{2.} Accéléré~: de $0.02 \div 0.1 = \qty{0.2}{m/s}$ à
\qty{1.0}{m/s}~; gain de \qty{0.2}{m/s} par intervalle ---
\qty{2}{m/s} par seconde.
\textbf{3.} Ralenti, et fortement. Le sable fournit une grande \omterm{def:g2:pushes-and-pulls:force}{force}
non compensée vers l'arrière --- un frottement profond --- et les
\omterm{def:g7:motion-average-speed:formula}{vitesses} dégringolent~: $1.0$, $0.7$, $0.45$, $0.25$,
\qty{0.1}{m/s}.
\textbf{4.} Trois intervalles de glissade uniforme à \qty{0.5}{m/s},
puis une accélération --- les écarts bondissant à $7.5$ puis
\qty{10.0}{cm}~: entre le troisième et le quatrième intervalle,
quelque chose s'est mis à pousser (une pente, une bourrade).
L'événement siège à la fin du troisième écart.
\textbf{5.} Les deux peuvent être vrais~: la moyenne mélange les
attentes et la route sur tout le trajet~; le radar lit l'instant du
carrefour. Le code de la route lit le radar --- les limitations
gouvernent la \omterm{def:g9:uniform-varied-motion:instant}{vitesse instantanée}, et aucune moyenne débonnaire
n'excuse un instant rapide.
\textbf{6.} Uniformément accéléré. Des paliers de \omterm{def:g7:motion-average-speed:formula}{vitesse} égaux en
des \omterm{def:g2:measuring-time:duration}{durées} égales désignent une poussée résultante constante et
invariable --- cause constante, croissance constante.
\textbf{7.} La poussée vers l'avant du sol sur les chaussures du
marcheur égale les traînées vers l'arrière de l'\omterm{def:g2:air-around-us:air}{air} et du sol sur le
corps~: compensées, d'où les \qty{0.8}{m/s} constants.
\textbf{8.} Croissance~: \qty{0.98}{m/s} par intervalle de
\qty{0.1}{s} --- \qty{9.8}{m/s} par seconde~: le $g$ local, reproduit
par la bande d'une pierre qui tombe.
\textbf{9.} L'affirmation tient~: \omterm{def:g7:motion-average-speed:formula}{vitesses} $1.2$, $1.15$, $1.10$,
$1.05$, \qty{1.0}{m/s} --- un ralentissement doux, du début à la fin.
\textbf{10.} Sortie $1.2 \times 3.6 = \qty{4.3}{km/h}$~; entrée
\qty{3.6}{km/h} --- l'une et l'autre bien en dessous de la référence
du cycliste~: imprudence acquittée~; la conduite, peut-être, est une
autre affaire.
\textbf{11.} «~Sur le plat, rien n'arrête une planche qui roule
sinon un frottement non compensé vers l'arrière --- les cours lisses
en fournissent peu, si bien que la planche garde l'essentiel de son
mouvement, exactement comme l'exige la loi des \omterm{def:g2:pushes-and-pulls:force}{forces}
compensées.~»
\textbf{12.} Par exemple~: «~Écarts égaux, allure tranquille --- nulle
\omterm{def:g2:pushes-and-pulls:force}{force} en trop ne trouble la file. Écarts qui changent, il faut une
cause~: une poussée non compensée, dit la clause.~»
\end{solution}
